\documentclass{article}

% Keep the official ICLR files in ../template/iclr2026/ while making this
% manuscript directly compilable from the main/ directory.
\makeatletter
\def\input@path{{../template/iclr2026/}}
\makeatother
\usepackage{iclr2026_conference,times}
\input{math_commands.tex}

\usepackage{fontspec}
\usepackage{xeCJK}
\setmainfont{TeX Gyre Termes}
\setCJKmainfont{AR PL UMing CN}
\setCJKsansfont{Droid Sans Fallback}
\setCJKmonofont{Droid Sans Fallback}
\XeTeXlinebreaklocale "zh"
\XeTeXlinebreakskip=0pt plus 1pt

\usepackage{amsmath,amssymb,mathtools}
\usepackage{booktabs,tabularx,threeparttable,multirow,array}
\usepackage{graphicx,xcolor,enumitem,float}
\usepackage{tikz}
\usetikzlibrary{arrows.meta,positioning,fit,calc,shapes.geometric,backgrounds}
\usepackage{pgfplots}
\pgfplotsset{compat=1.18}
\usepackage{hyperref}
\usepackage{url}

\hypersetup{colorlinks=true,linkcolor=blue,citecolor=blue,urlcolor=blue}
\setlist{nosep,leftmargin=*}
\newcolumntype{Y}{>{\centering\arraybackslash}X}
\newcommand{\scc}{\textsc{SC-CL}}
\newcommand{\bone}{\textsc{B1}}
\newcommand{\preset}{\textcolor{gray}{\textsuperscript{p}}}
\newcommand{\todoresult}{\textcolor{gray}{待运行验证}}

\title{开放世界知识图谱关系候选排序的可审计结构反事实对比学习}
\author{}

% \iclrfinalcopy % Keep commented for the anonymous review template.
\begin{document}

\maketitle

\begin{abstract}
生成式知识图谱关系预测可用候选级对比学习强化关系辨别，但困难候选的来源与真假边界常不明确；在开放世界知识图谱中，未标注关系也不能直接视为假事实。本文提出结构反事实对比学习（Structure-Counterfactual Contrastive Learning, \scc）：从训练关系词表中过滤正确关系及同一实体对的训练已知关系，结合关系结构画像与冻结 GRIP 基线（\bone）的分数违例度，按带均匀下限的分布无放回抽取九个候选。固定候选库接入 GRIP Stage~2 的候选级 InfoNCE，原始生成损失保持不变。NELL23K 的一次 seed=2026 运行在 4,943 条测试问题上取得候选重排 EM/MRR 91.52\%/95.27\%；对相同的十选一题目直接解码，\scc{} 与 \bone{} 的 EM 分别为 87.03\% 和 87.46\%。同候选的前 200 条问题上，两者的重排 MRR 分别为 94.96\% 和 93.36\%。这些结果证明候选构造与训练链路可复核，并给出有限的排序观察；全量 \bone{} 重排、同预算随机对照和多种子仍是判断方法收益所需的证据。
\end{abstract}

\section{引言}

给定实体对及其图上下文，关系预测需要从关系词表中识别相应的关系。GRIP 将图结构融入语言模型参数，通过生成式问答完成这一任务 \citep{feng2025grip}。在原有生成目标之外加入候选级对比损失，可以让目标关系与替代关系直接竞争；这一训练信号的作用取决于替代关系如何选取。

候选过于容易时，提供的排序边界信号可能有限；只按模型高分挑选时，候选又可能反映偶然的分数混淆。更关键的是，知识图谱通常不完整：$(h,r,t)$ 未被标注，并不等于该事实为假。困难候选因此同时面临\emph{信息量}和\emph{事实不确定性}两项要求。本文聚焦一个可检验的问题：能否在保持 GRIP 原始生成路径的同时，构造来源可追溯、排除训练已知关系的候选库，并用它训练关系排序边界？

为此，我们提出结构反事实对比学习（\scc）。首先，在训练关系词表内排除正关系及同一实体对的 train-known 关系；随后，用训练图的关系画像衡量结构相似度，用冻结的 generation-only 基线 \bone{} 衡量候选分数违例度。两种信号经带均匀下限的分布组合，每题无放回采样九个候选，并在 Stage~2 前冻结。训练仅为原生成目标增加候选级 InfoNCE，不改变 Stage~1、原始提示或自由生成评估。这里的“结构反事实”指对已观测实体对替换关系标签，不表示剩余候选都是真实负例。图~\ref{fig:mechanism} 展示训练流程。

评估在相同的十选一题目上分别考察模型直接解码与逐候选评分。NELL23K 的一次 seed=2026 运行完成了候选库、Stage~2 训练和全测试重排，得到候选 EM/MRR 91.52\%/95.27\%；这一绝对值缺少全量 \bone{} 重排对照。在目前相同候选列表的前 200 条问题上，\scc{} 与 \bone{} 的 MRR 分别为 94.96\% 和 93.36\%；全量直接解码 EM 分别为 87.03\% 和 87.46\%。前缀子集、训练预算差异与单种子设计都限制了因果解释。本文因此报告可复核的候选构造协议和有限配对观察，将组合信号的独立收益留给同预算随机对照与多种子实验检验。

本文的具体贡献是：定义同时利用结构相似度和冻结模型违例度的训练期候选分布；将 train-known 过滤、固定候选库和来源记录落实为可审计约束；在同一知识图谱上分别报告候选重排与直接解码结果，明确两种答案选择规则对应的指标。全文限于固定实体对的关系预测，不从当前观察外推跨数据集收益。

\section{相关工作}

\paragraph{关系预测与生成式图推理。}
TransE、ComplEx、RotatE 和 ConvE 通过不同的三元组打分函数建模实体与关系 \citep{bordes2013transe,trouillon2016complex,sun2019rotate,dettmers2018conve}。KG-BERT 将三元组判断转为文本建模 \citep{yao2019kgbert}，SimKGC 则在预训练语言模型上使用多种实体负例进行对比学习 \citep{wang2022simkgc}。这些工作通常评估缺失实体或三元组判断，不直接给出 GRIP 中生成损失与候选关系损失如何并存的训练协议。GRIP 以图上下文和生成式问答学习参数化图推理 \citep{feng2025grip}；本文固定其 Stage~1、提示和生成路径，改变的是 Stage~2 的候选关系来源。

\paragraph{困难候选与假阴性。}
结构感知负采样 SANS 和对抗式采样 KBGAN 分别从图结构与模型反馈寻找困难负例 \citep{ahrabian2020sans,cai2018kbgan}。对比学习中的难例还可能与目标语义相近；去偏对比学习和 ProGCL 因而指出采样难度与假阴性风险不能混为一谈 \citep{chuang2020dcl,xia2022progcl}。本文从这一张力出发，将训练图关系画像与冻结基线的分数违例度组合，但只把 train-known 过滤视为对已知冲突的排除，不把其余候选宣称为经验证的假事实。带均匀下限的采样分布和冻结 bank 进一步限定了候选的来源与更新时间。

\paragraph{评估边界。}
传统 link prediction 常在实体替换候选上报告 MRR/Hits，CoDEx 还提供经过人工核验的实体替换困难负例 \citep{safavi2020codex}。本文固定实体对、替换关系标签；两种候选空间不能直接共享负例真值。我们分别报告直接解码和给定候选排序，并将方法主张限制在具有相同候选列表与可核对训练记录的比较上。当前只有 NELL23K 的 SC-CL 运行具备正文主表所需的完整输出；WN18RR 的单次直接解码结果列于附录，尚不能作为跨图收益证据。

\section{问题定义}

\subsection{训练图与候选关系}

令训练知识图谱为 $\mathcal{G}_{\mathrm{tr}}=(\mathcal{V},\mathcal{E}_{\mathrm{tr}})$，其中
$\mathcal{E}_{\mathrm{tr}}\subseteq \mathcal{V}\times\mathcal{R}_{\mathrm{tr}}\times\mathcal{V}$，
$\mathcal{R}_{\mathrm{tr}}$ 是训练关系词表。一个关系问答样本表示为
$q=(h,t,r^+)$，其中 $r^+$ 是目标关系，$h,t$ 是有序实体对。定义同一实体对在训练图中的已知关系集合为
\begin{equation}
    T_{\mathrm{tr}}(h,t)=\{r:(h,r,t)\in\mathcal{E}_{\mathrm{tr}}\}.
\end{equation}
候选池为
\begin{equation}
    P_q=\mathcal{R}_{\mathrm{tr}}\setminus\bigl(\{r^+\}\cup T_{\mathrm{tr}}(h,t)\bigr).
    \label{eq:pool}
\end{equation}
式~\ref{eq:pool} 只保证候选没有作为该实体对的\emph{已知训练关系}出现；它不声称候选事实在真实世界中为假。这一点是本文“反事实”术语的边界，也是开放世界知识图谱中避免过度解释的必要条件 \citep{ji2022survey,hogan2021kg}。

\subsection{两个评估目标}

本文区分两种答案选择规则。第一，模型在列出十个关系选项的题目上直接解码关系文本，评价 Exact Match（EM）和 OOV；它衡量模型是否保持原始 GRIP 的生成路径。第二，评估程序给定同一题目的 $\{r^+,r_1^-,\ldots,r_K^-\}$，用关系 continuation score $z_\theta(q,r)$ 逐项排序，评价 MRR、Hits@1/3/10 和 wrong-in-list。后者更直接测量候选 InfoNCE 是否改变关系边界，但不是模型直接解码所得的准确率。

给定冻结 bank $B(q)=\{r_1^-,\ldots,r_K^-\}$，候选损失写为
\begin{equation}
    \mathcal{L}_{\mathrm{cand}}(q)=-\log
    \frac{\exp(z_\theta(q,r^+)/\tau_c)}
    {\exp(z_\theta(q,r^+)/\tau_c)+\sum_{k=1}^K\exp(z_\theta(q,r_k^-)/\tau_c)}.
    \label{eq:infonce}
\end{equation}
其中 $\tau_c$ 是候选温度。训练阶段保留原始生成损失 $\mathcal{L}_{\mathrm{gen}}$，只加入
\begin{equation}
    \mathcal{L}=\mathcal{L}_{\mathrm{gen}}+\lambda_{\mathrm{cand}}\mathcal{L}_{\mathrm{cand}}.
    \label{eq:total}
\end{equation}
Stage 1 图上下文训练和 B1 generation-only 路径不变，使得匹配预算比较具有清晰的归因边界。

\section{结构反事实对比学习}

图~\ref{fig:mechanism} 概括 SC-CL 从训练期候选构造到 Stage~2 联合优化的流程：训练图和训练问题给出可用关系与结构画像，冻结 B1 的训练集分数提供模型违例度；两种信号在已过滤的候选池中混合采样，形成固定的九候选 bank；最后，GRIP Stage~2 联合优化候选损失与原始生成损失。训练期间 bank 只被读取，不根据 valid/test 预测重新采样。自由生成和候选排序的评估目标按上一节分别定义。

\begin{figure}[H]
    \centering
    % Extracted, without resampling, from ../artifacts/结构引导与候选级对比优化流程图.svg.
    \includegraphics[width=\linewidth]{figures/sccl_structure_guided_svg_content.png}
    \caption{\scc{} 的训练流程。A：训练图关系画像与冻结 \bone{} 分数提供双信号，经过 train-known 过滤与无放回采样，形成每题 $K=9$ 的固定候选库。B：Stage~2 保留原始生成损失，并加入候选级 InfoNCE 共同优化 LoRA。图中损失求和省略了候选权重；严格定义见式~\ref{eq:total}。}
    \label{fig:mechanism}
\end{figure}

图中顶部 $C_0$ 的加载路径为概念性示意：当前实现的 Stage~2 从 Stage~1 adapter 初始化，冻结的 \bone{} 仅用于离线候选评分，不参与训练时的梯度更新。

\subsection{训练图关系画像}

对每一个 $r\in\mathcal{R}_{\mathrm{tr}}$，我们从 $\mathcal{E}_{\mathrm{tr}}$ 计算五维关系画像
\begin{equation}
    \phi(r)=\bigl(\log(1+|E_r|),\log(1+|H_r|),\log(1+|T_r|),
    \overline{d}_{H,r},\overline{d}_{T,r}\bigr),
    \label{eq:profile}
\end{equation}
其中 $|E_r|$ 为关系边数，$|H_r|$ 和 $|T_r|$ 分别为关系覆盖的头实体数与尾实体数，
$\overline{d}_{H,r}$ 是关系边上头实体的平均出度，$\overline{d}_{T,r}$ 是尾实体的平均入度。前两类计数使用 $\log(1+x)$，避免高频关系完全支配距离；度统计保留图的局部连接性。

对目标关系 $r^+$ 和候选 $r$，定义归一化结构相似度
\begin{equation}
    s_{\mathrm{str}}(q,r)=1-\frac{1}{5}\sum_{j=1}^{5}
    \frac{|\phi_j(r^+)-\phi_j(r)|}
    {|\phi_j(r^+)|+|\phi_j(r)|+\epsilon},
    \label{eq:similarity}
\end{equation}
其中 $\epsilon$ 防止两个稀疏特征同时为零。实现将该值裁剪到 $[0,1]$。这个分数只表达“结构画像相近”，不表达两个关系的语义等价性，也不替代对当前实体对的已知关系过滤。

\subsection{冻结 B1 的边界违例度}

\bone{} 是原始 GRIP 的 generation-only 基线。我们在训练问题上离线计算全关系词表的 continuation score。若离线记录已有 pairwise confusion，则直接使用裁剪到 $[0,1]$ 的 confusion；否则对候选和正关系的分数差使用平滑违例度
\begin{equation}
    V_{\bone}(q,r)=\sigma\left(\frac{z_{\bone}(q,r)-z_{\bone}(q,r^+)}{\tau_v}\right).
    \label{eq:violation}
\end{equation}
候选越接近或超过正关系，$V_{\bone}$ 越大。B1 的输出在 bank 构造后冻结，因而 SC-CL 不会把自身训练后的 test 预测反过来用于采样。

\subsection{结构反事实候选分布}

将结构相似度和分数违例度合成为
\begin{equation}
    u(q,r)=\alpha s_{\mathrm{str}}(q,r)+\beta V_{\bone}(q,r),
    \qquad r\in P_q.
    \label{eq:utility}
\end{equation}
候选分布使用均匀下限与 utility softmax 的混合：
\begin{equation}
    p(r\mid q)=\frac{\rho}{|P_q|}+(1-\rho)
    \frac{\exp(u(q,r)/\tau_s)}{\sum_{r'\in P_q}\exp(u(q,r')/\tau_s)}.
    \label{eq:sampling}
\end{equation}
我们从 $p(r\mid q)$ 无放回抽取 $K=9$ 个负关系。默认参数为 $\rho=1/3$、$\alpha=1$、$\beta=1$、$\tau_s=0.25$。如果 $|P_q|<K$，样本进入失败清单而不是静默重复候选；正式实验必须报告该清单及其原因。

这里的“反事实”可写成一个受限操作：在已观测的 $(h,t,r^+)$ 上将标签替换为 $r\in P_q$。$r$ 不是 $T_{\mathrm{tr}}(h,t)$ 中的关系，因此该替换没有复现训练图已知事实；但它仍可能对应真实世界中的未标注事实。这个定义保留了知识图谱的开放世界语义，并让“合法负例”成为数据协议而非未经证实的真值断言。

\subsection{冻结 bank 与候选学习}

SC-CL 的预处理阶段写入每个问题的 question ID、正关系、九个候选、每个候选的结构相似度、违例度、utility、B1 score gap、采样参数、split 和 known-pair 过滤记录。bank manifest 至少满足以下门：

\begin{enumerate}
    \item 所有问题的 \texttt{split} 必须为 \texttt{train}，\texttt{filter\_splits} 只能包含 \texttt{train}；
    \item 每个样本恰有一个正关系和 $K=9$ 个互异候选，候选均属于 $\mathcal{R}_{\mathrm{tr}}$；
    \item 候选与 $T_{\mathrm{tr}}(h,t)$ 不相交，且 question ID 与完整 task JSON 一一对齐；
    \item bank hash、脚本版本、随机种子和结构采样参数被写入 manifest，训练后不得覆盖原文件。
\end{enumerate}

训练时，候选 continuation score 代入式~\ref{eq:infonce}，生成损失仍按原始 GRIP 路径计算。该设计将 SC-CL 的影响集中到“给定候选集合如何区分关系”，而不是同时改变上下文构造或自由解码。

\subsection{复杂度与实现边界}

关系画像只需一次线性扫描训练边，复杂度为 $O(|\mathcal{E}_{\mathrm{tr}}|)$；每个问题的候选 utility 在关系词表上计算，复杂度约为 $O(|\mathcal{R}_{\mathrm{tr}}|)$，采样后 bank 大小为 $O(K|Q_{\mathrm{tr}}|)$。结构画像与采样本身不需要 GPU；冻结 B1 的全关系 continuation scoring 和后续候选训练则需要额外模型计算，实测成本见第 5 节。确定性、过滤和非训练 split 拒绝通过实现测试检查。

\section{实验与结果}

\subsection{任务、评估方式与证据范围}

本文评价固定实体对上的十选一关系预测。两种评估使用含相同候选选项的题目：直接解码时，模型自行输出关系文本，以 Exact Match（EM）衡量；候选重排时，评估程序逐一代入十个选项，计算 teacher-forced continuation score，再以最高分关系计算 EM、MRR 和 Hits。重排 EM 等于 H@1。选项相同并不意味着答案选择规则相同，两套 EM 不能直接相减来计算训练收益，也不能把重排 EM 当作原始 GRIP 的生成式准确率。

当前具备训练、直接解码和候选重排完整输出、可进入正文主表的 SC-CL 运行位于 NELL23K，使用 Qwen2.5-7B-Instruct 和 seed=2026；测试集包含 4,943 条关系问题。FB15K237-2 的 B1 与历史 InfoNCE-i0 各有 10,000 条候选评分，但 i0 的候选文件、LoRA rank、训练轮数和损失配置不同，不能构成 SC-CL 的跨数据集复现；该结果仅列于附录表~\ref{tab:historical-fb}。WN18RR 已有单种子探索性直接解码测试，但尚无完整候选重排和匹配对照；CoDEx-M 尚无方法训练结果。补充状态见附录~\ref{sec:wn-exploratory}，后续实验范围见附录~\ref{sec:future-protocol}。

\subsection{先核对候选库与训练链路}

NELL23K 的候选库由训练问题构造，写入 3,253 行，每行九个互异关系，\texttt{rows\_failed=0}。冻结 bank 的 wiring 报告确认训练来源、任务问题对应关系和 manifest 记录；SC-CL Stage 2 完成 10 个 epoch、240 个优化步，保存了 adapter、自由生成预测和候选重排输出。这些检查支持\emph{该次运行的来源与执行完整性}，不证明剩余候选均为真实负例。尤其是与持出正例的重叠和开放世界未知事实仍需单独审计。

\subsection{候选排序：有限配对观察与全量结果}

表~\ref{tab:nell-results} 将严格配对的前 200 条问题与只有 SC-CL 的全量闭集结果分开。前 200 条使用相同问题和相同候选列表；B1 的完整测试候选分数尚未完成，故全量一行没有相应的 B1 对照。这样的排列让可比较的差异与单模型绝对水平各自承担不同的证据作用。

\begin{table}[t]
\caption{NELL23K 已观察结果（百分比；单次 seed=2026）。前 200 条为相同问题和候选列表上的 B1/SC-CL 配对比较；full 为仅有 SC-CL 的 4,943 条测试问题。直接解码和候选重排使用相同的十选一题目，但答案选择规则不同。}
\label{tab:nell-results}
\centering\scriptsize
\begin{tabular}{llrrrrr}
\toprule
评估 & 模型 & $n$ & EM/H@1 & MRR & H@3 & H@10 \\
\midrule
闭集 prefix-200 & \bone & 200 & 88.50 & 93.36 & 97.50 & 100.00 \\
闭集 prefix-200 & \scc-9 & 200 & 91.00 & 94.96 & 99.50 & 100.00 \\
闭集 full & \scc-9 & 4,943 & 91.52 & 95.27 & 99.07 & 100.00 \\
\midrule
直接解码 full & \bone & 4,943 & 87.46 & -- & -- & -- \\
直接解码 full & \scc-9 & 4,943 & 87.03 & -- & -- & -- \\
\bottomrule
\end{tabular}
\end{table}

在配对子集上，SC-CL 的 MRR 比 B1 高 1.60 个百分点，候选 EM 高 2.50 个百分点；逐题 top-1 变化为 9 条由错转对、4 条由对转错、173 条共同正确和 14 条共同错误。净变化为五题，同时存在收益与回退。这里使用的是\emph{前} 200 条，而非随机代表样本；单种子结果也没有不确定性估计，因此不能把差异外推到全部测试问题。全量 SC-CL 闭集 MRR 为 95.27\%，描述的是该模型的绝对排序水平；缺少全量 B1 候选评分时，无法计算对应的全量训练增益。

\subsection{直接解码结果与解释边界}

同一批十选一测试问题上，B1 与 SC-CL 的直接解码 EM 分别为 87.46\% 和 87.03\%，描述性差异为 $-0.43$ 个百分点。当前数据没有显示生成能力提升；这一下降是否稳定或是否来自候选训练，尚需同预算、多种子对照检验。SC-CL 的重排 EM 91.52\% 高于自身直接解码 EM 87.03\%，只说明逐候选评分改变了答案选择方式，不能当作训练带来的 4.49 个百分点收益。GRIP 原论文在 NELL23K 上报告 87.74\% 的生成式准确率，但其附录记录 4,952 条测试关系，本地为 4,943 条；该数值只作外部背景，不参与主表的配对差异计算 \citep{feng2025grip}。

\subsection{计算成本与证据缺口}

该次 NELL 运行在训练集上为 644,094 个关系候选生成冻结 B1 分数，用时 19,027 秒，记录的峰值显存为 15.21 GiB。SC-CL Stage 2 处理 12,148,720 个输入 token，用时约 80,870 秒；全测试候选重排用时 3,567 秒。这些是单次运行的绝对成本。由于缺少同硬件、同预算的 Random-9 和 B1 完整训练/评估对照，当前不能计算可靠的单位收益成本比。

全量 B1 闭集评分、同预算 Random-9 和频率匹配随机对照仍缺失，因此目前无法分别估计候选损失与候选构造的作用。持出正例重叠和开放世界未判定也尚未完成完整审计。附录~\ref{sec:future-protocol} 列出后续比较的判定规则；本节的实证范围限于一次可复核运行与前 200 条的描述性配对观察。

\section{讨论与局限}

\paragraph{候选排序差异的解释。}
同候选的前 200 条问题上，SC-CL 的 MRR 高于 B1，这与困难候选可能提供更清晰排序信号的设想相容。但这一比较同时改变了候选损失和训练预算，也没有把结构画像与 B1 违例度单独剥离。前缀子集的分布可能不同于剩余测试问题，且全量 B1 闭集分数缺失。因此，目前只能报告配对观察，不能把差异归因于组合采样机制。

\paragraph{排序与生成的取舍。}
SC-CL 在逐候选评分下的 MRR 为 95.27\%，但直接解码 EM 为 87.03\%，略低于 B1 的 87.46\%。两种评估的题目都列出十个选项；差别在于模型自行解码，还是评估程序逐一代入选项评分。因此，两个数值不能互为训练收益证据。如果同预算复现仍呈现排序改善而直接解码不增，方法的用途应限定为候选重排；原始 GRIP 生成任务的整体收益需要单独证实。

\paragraph{过滤的边界。}
train-known 过滤仅排除训练图中该实体对已有的关系。剩余候选可能在持出 split 或真实世界中成立，因此“结构反事实”是训练标签替换的操作名，不是经过验证的假事实。持出正例重叠与开放世界未判定需要分开统计；CoDEx 的实体替换硬负例也不能为本文的关系替换候选提供真值 \citep{safavi2020codex}。若后续审计发现高于随机候选的冲突风险，应相应收窄方法主张或调整训练设计。

\paragraph{成本与适用范围。}
冻结 bank 使候选来源可追溯，但需要离线计算 B1 的全关系分数，也不能随训练中的模型更新。现有资源记录只覆盖 NELL23K 的一次运行，尚不能判断收益是否抵消额外成本。方法针对固定实体对上的关系预测；图级组合问答需要重新定义反事实单位和评估目标。跨图适用性仍取决于 FB15K237-2、WN18RR 和 CoDEx-M 上的同契约评估。

\section{结论}

本文提出 SC-CL，在训练关系词表内过滤同一实体对的已知关系，再以结构画像和冻结基线的分数违例度构造九候选库，用于 GRIP Stage~2 的候选级 InfoNCE。固定候选库及其来源记录使训练流程可复核；未观测关系仍保留开放世界中的真值不确定性。

NELL23K 的一次运行在 4,943 条问题的候选重排上得到 91.52\% EM 和 95.27\% MRR；同候选的前 200 条问题中，SC-CL 与 B1 的 MRR 分别为 94.96\% 和 93.36\%。与此同时，十选一题目上的直接解码 EM 分别为 87.03\% 和 87.46\%。当前证据表明方法能够执行且值得检验其排序作用，但尚不能证明组合信号的独立收益、生成改善或跨数据集稳定性。同预算候选对照、多种子复现和开放世界冲突审计是上述主张成立的必要条件。


\clearpage
\appendix
\section{补充结果与审计协议}

\subsection{历史 FB 候选评分的来源边界}

\begin{table}[h]
\caption{FB15K237-2 历史候选评分（百分比，10,000 条测试问题）。InfoNCE-i0 使用历史候选文件和不同训练配置，不是 SC-CL 的跨数据集结果。}
\label{tab:historical-fb}
\centering\scriptsize
\begin{tabular}{lrrrrr}
\toprule
方法 & 自由 EM & 候选 EM/H@1 & MRR & H@3 & H@10 \\
\midrule
\bone & 73.23 & 79.91 & 88.01 & 95.80 & 100.00 \\
历史 InfoNCE-i0 & 88.69 & 91.62 & 95.41 & 99.20 & 100.00 \\
\bottomrule
\end{tabular}
\end{table}

两组评分均来自相同的本地 processed test，哈希为
\texttt{c77dde51aafc7b1ce3a5011ade914343c2689f8b38b12774ea4840ec1c9068ed}。
历史 i0 使用 \texttt{fb15k237\_i0\_paired\_candidates.json}，LoRA rank 为 24、Stage 2
训练 5 个 epoch、候选损失权重为 0.1；这些条件与本文 SC-CL 的训练契约不一致。

\subsection{WN18RR 单种子探索性生成评估}
\label{sec:wn-exploratory}

WN18RR 的训练期 SC-CL bank 覆盖 2,709 条关系问答，每题九个候选且无失败行；wiring 记录为 train-only，并与任务问题对齐。使用 seed=2026 时，在完整验证集上从候选损失权重中选择 $\lambda_{\mathrm{cand}}=0.25$，随后只对获选模型执行一次 3,134 条测试问题的直接解码评估，得到 EM 86.89\%。本地 B1 在同一测试文件上的 EM 为 32.67\%。原始 GRIP 论文在 WN18RR 报告 64.30\%，但这是五次运行均值；本地 B1 与该值的差距尚未解释，也没有完成 WN 的同候选重排和匹配随机对照。因此，这里仅记录可复核的探索性生成结果，不把外部数字或本地差异写成 SC-CL 的稳健跨图收益 \citep{feng2025grip}。

原始预测与配置分别位于 \texttt{results/hpo/wn18rr\_sccl\_seed2026/lambda0p25\_final\_test/} 和 \texttt{results/runs/wn18rr\_sccl\_hpo\_s2026\_lc0p25/}；这两个目录相对于实验模块 \texttt{08\_experiments/Hard\_Negative\_Contrastive\_GRIP/}。完整评估契约仍要求候选重排、生成输出、配置哈希及匹配控制共同齐备。

\subsection{后续比较与可证伪规则}
\label{sec:future-protocol}

完整主比较需要相同基座、Stage 1 adapter、任务文件、提示、候选数、候选损失设置和
Stage 2 预算。\bone{} 是原始 generation-only 锚点；Random-9 衡量加入候选损失后的
差异；频率匹配随机检验结构画像是否只是关系频率代理；Profile-only 与 Violation-only
分别检验两个信号，ErrorMemory-9 仅作辅助对照。若 SC-CL 不优于频率匹配随机，组合
采样优势的主张必须撤回；若自由生成 EM 或 OOV 恶化，需单独报告生成副作用。

计划在 NELL23K 和 FB15K237-2 上使用训练种子 2026--2030，报告逐种子结果与问题级
配对区间；$K\in\{3,9,15\}$、$\rho\in\{0,1/3,2/3\}$ 和
$\tau_s\in\{0.10,0.25,0.50\}$ 用于敏感性检查。WN18RR 仍需补齐同候选重排、匹配控制和基线复现解释，CoDEx-M 则需完成任务文本、
adapter、bank 和完整评估的对齐审计后才能进入迁移比较。当前 CoDEx-M 已转换官方正例
split 和固定实体对的关系候选，但 Stage 1/Stage 2 推理及方法训练尚未完成；其通过替换
实体构造的官方硬负例不作为关系候选的真值标签。

每个正式运行需记录 bank 的训练来源、question ID 对齐、每行九个互异候选、train-known
过滤结果、持出正例重叠、数据和 adapter 哈希，以及候选评分与自由生成原始输出。
未在训练图出现且未被其他证据确认的候选列为“开放世界未判定”，不计入已证实假阴性。
任何 split 泄漏或任务对齐失败均使该运行退出主比较。

\subsection{默认超参数}

\begin{table}[h]
\caption{SC-CL 的默认设置和正式运行前需要登记的字段。}
\label{tab:hyper}
\centering\small
\begin{tabular}{ll}
\toprule
字段 & 默认值 \\
\midrule
候选数 $K$ & 9 \\
均匀混合系数 $\rho$ & $1/3$ \\
结构权重 $\alpha$ & 1.0 \\
B1 违例权重 $\beta$ & 1.0 \\
采样温度 $\tau_s$ & 0.25 \\
候选损失温度 $\tau_c$ & 运行配置登记 \\
候选损失权重 $\lambda_{\mathrm{cand}}$ & 1.0 \\
bank 随机种子 & 2026 \\
训练种子 & 2026, 2027, 2028, 2029, 2030 \\
过滤来源 & train-only \\
\bottomrule
\end{tabular}
\end{table}

\subsection{候选 bank 的最小 schema}

每一行 bank 至少包含：\texttt{question\_id}、\texttt{split}、\texttt{head}、\texttt{tail}、\texttt{gold\_relation}、
\texttt{negative\_relations}、\texttt{negative\_source}、\texttt{filter\_splits}、\texttt{counterfactual\_operation}、
\texttt{counterfactual\_validation}、\texttt{K}、\texttt{uniform\_mix}、\texttt{structure\_weight}、
\texttt{violation\_weight}、\texttt{temperature}、\texttt{seed} 和 sampler 版本。每个候选记录还应包含
\texttt{structure\_similarity}、\texttt{score\_violation}、\texttt{utility}、\texttt{candidate\_score}、
\texttt{gold\_score}、\texttt{score\_gap} 和 \texttt{counterfactual\_valid}。字段名以项目的实际 JSONL
实现为准，论文只要求这些语义信息可审计。

\subsection{审计伪代码}

\small
\begin{verbatim}
for q in train_questions:
    assert q.split == "train"
    P = train_relations - {q.gold} - known_train_relations(q.head, q.tail)
    for r in P:
        s = profile_similarity(phi[q.gold], phi[r])
        v = b1_score_violation(q, r)
        u[r] = alpha * s + beta * v
    negatives = sample_without_replacement(
        P, mixture_softmax(u, rho, tau), K=9)
    write_frozen_bank(q, negatives, manifest=provenance)
\end{verbatim}

该伪代码强调三个不可省略的顺序：先限定训练 split，再过滤实体对已知关系，最后才计算采样分布。正式 wiring 检查还要确认训练 task 的 question ID、正关系文本和候选关系文本保持一致。

\subsection{失败分类}

\begin{table}[h]
\caption{建议的失败分析标签。}
\label{tab:failure}
\centering\scriptsize
\begin{tabular}{p{0.23\linewidth}p{0.65\linewidth}}
\toprule
标签 & 判定规则 \\
\midrule
结构稀疏 & 目标关系或候选关系的画像特征大部分为零，profile 相似度不稳定。 \\
已知关系冲突 & 候选属于当前实体对 train-known 关系，说明过滤或 bank 版本错误。 \\
开放世界不确定 & 候选未在训练图出现，但在外部/测试事实中可能成立，不能当作确定假阴性。 \\
B1 高置信混淆 & 候选 score 高于 gold，且自由生成与闭集排序结论不一致。 \\
候选不足 & 过滤后 $|P_q|<9$，必须进入失败清单，不能重复填充。 \\
生成-排序分歧 & 闭集 MRR 改善但自由生成 EM 下降或不变。 \\
\bottomrule
\end{tabular}
\end{table}

\subsection{伦理与数据说明}

本文使用本地知识图谱基准和公开模型技术报告作为背景，不新增个人敏感数据。当前正文
只把带有原始预测、配置、manifest 和 hash 的 NELL 单种子结果作为实测；历史 FB
InfoNCE-i0 另行标注，未完成协议不用于性能承诺。正式提交前应补充数据许可、模型许可
证、代码版本和可公开的 bank/预测文件说明。



\bibliographystyle{iclr2026_conference}
\bibliography{references}

\end{document}
